
\section{Reflection positivity and a H\"older inequality for signed laws}
\label{sec:positivity}

\paragraph{What this section proves and how to read it.}
This section proves that a signed $72$-fold average can still be
bounded below by $\alpha^{72}$ and above by diagonal averages,
provided the averaging law is positive across mirror symmetries.
Precisely, it defines cuts and reflection positivity, proves that
domination by a positive $2\times2$ block at one cut permits
Cauchy--Schwarz for signed laws, iterates this over seven cuts to
obtain the signed H\"older inequality, and constructs the $72$ slots
with their seven cuts and six queries. Signed laws appear because
the comparison law $\nu$ of Part~III cannot be reflection positive;
the block domination is what replaces positivity for $\nu$. Inputs:
only finite linear algebra. Outputs: \cref{thm:holder} (hence
$\alpha^{72}\le Y$) and the block-domination calculus used in
Part~III. A first reading should follow the definitions,
\cref{thm:holder}, and the slot construction in
\cref{subsec:slot-array}; the detailed one-cut lemmas can be treated
as reference. Until \cref{subsec:slot-array}, $V$ is an arbitrary
finite slot set; from there on $V=(\Z/6\Z)^2\times\{0,1\}$.

\paragraph{Notation.}
Throughout, $V$ is a finite set of \emph{slots}, $d=|V|$, and $\mathcal X$ is
a finite set of labels. For a finite signed measure $\nu$ on $\mathcal X^V$ and
real functions $g_v$ on $\mathcal X$ we write
\[
 \Lambda_\nu\big((g_v)_{v\in V}\big)=\int\prod_{v\in V}g_v(z_v)\,d\nu(z),
 \qquad
 \Lambda_\nu(g)=\Lambda_\nu\big((g)_{v\in V}\big).
\]
A map $h\colon V\to\R^{\mathcal X}$ is an \emph{assignment}, and
$\Lambda_\nu(h)=\Lambda_\nu((h(v))_{v\in V})$. For a map $\Psi$ between finite
sets, $\Psi_*\kappa$ is the image of a measure $\kappa$; for a permutation
$\sigma$ of $V$ we write $\sigma_*\kappa$ for the image of $\kappa$ under the
map $z\mapsto z\circ\sigma=(z_{\sigma v})_{v\in V}$ of $\mathcal X^V$.

\subsection{Cuts and cut matrices}\label{subsec:pos:cuts}

Laws obtained by restricting a law to an event, by repair, or as
comparison laws are in general not reflection positive. Two facts replace positivity at a single cut: measures that the cut
fixes samplewise are ordered in the Loewner order at that cut
(\cref{lem:commoncut}), and a measure bracketed at one cut by a reflection
positive law inherits that law's mixed H\"older bound
(\cref{lem:record:block-gram}, \cref{rem:cc:mixed}). This subsection fixes the
language.

\begin{definition}[Cuts and reflection positivity]\label{def:rp}
A \emph{cut} of $V$ is a fixed-point-free involution $\vartheta$ of $V$ together
with a decomposition $V=V_+\sqcup V_-$ with $V_-=\vartheta V_+$; $V_+$ is its
\emph{plus half}. We name a cut by its involution when the plus half is
understood. A finite signed measure $\nu$ on $\mathcal X^V$ is
\emph{reflection positive} (RP) at $\vartheta$ if it is invariant under
$(z_v)_v\mapsto(z_{\vartheta v})_v$ and
\[
 \int F\big(z_{V_+}\big)\,F\big((z_{\vartheta u})_{u\in V_+}\big)\,d\nu(z)\ \ge\ 0
 \qquad\text{for every }F\colon\mathcal X^{V_+}\to\R .
\]
Equivalently (\cref{def:cc:order}), its cut matrix
$\big(\nu(z_{V_+}=x,\ z_{\vartheta V_+}=y)\big)_{x,y\in\mathcal X^{V_+}}$ is
symmetric and positive semidefinite.
\end{definition}

Positive multiples and sums of measures that are RP at $\vartheta$ are RP at
$\vartheta$, and so are tensor products of such measures \eqref{eq:cc:kron}.
Since $|V|=2|V_+|$, a slot set with a cut has an even number of slots.

\begin{definition}[Cut matrix, cut order and cut form]\label{def:cc:order}
Fix a cut $\vartheta$ with halves $V_\pm$. For a finite signed measure $\nu$
on $\mathcal X^V$ its \emph{cut matrix} at $\vartheta$ is
\[
 C_\vartheta(\nu)_{x,y}=\nu\big(z_{V_+}=x,\ z_{\vartheta V_+}=y\big)
 \qquad(x,y\in\mathcal X^{V_+}),
\]
where $z_{V_+}=x$ means $z_u=x_u$ and $z_{\vartheta V_+}=y$ means
$z_{\vartheta u}=y_u$ for all $u\in V_+$. For real symmetric matrices,
$A\preceq B$ means that $B-A$ is positive semidefinite (the Loewner order of
\cref{def:prelim:operators}). For finite signed
measures $\kappa,\kappa'$ on $\mathcal X^V$ we write $\kappa\succeq_\vartheta\kappa'$,
or $\kappa'\preceq_\vartheta\kappa$, if $C_\vartheta(\kappa-\kappa')$ is
symmetric and positive semidefinite; thus $\nu$ is RP at $\vartheta$ if and
only if $\nu\succeq_\vartheta0$. The \emph{cut form} of $\kappa$ at $\vartheta$
is the bilinear form
\[
 Q_\kappa(X,Y)=\ip{X}{C_\vartheta(\kappa)Y}
 =\int X\big(z_{V_+}\big)\,Y\big((z_{\vartheta u})_{u\in V_+}\big)\,d\kappa(z)
\]
on real functions $X,Y$ on $\mathcal X^{V_+}$ (\emph{half-functions}), with the
Euclidean inner product of $\R^{\mathcal X^{V_+}}$.
\end{definition}

The map $z\mapsto\big(z_{V_+},(z_{\vartheta u})_{u\in V_+}\big)$ is a bijection
of $\mathcal X^V$ onto $\mathcal X^{V_+}\times\mathcal X^{V_+}$, so a measure is
determined by its cut matrix at any one cut. The cut form $Q_\kappa(X,Y)$ is
bilinear in $(X,Y)$ and linear in $\kappa$; if $\kappa\succeq_\vartheta0$ it is
a symmetric positive semidefinite form, and $\kappa\preceq_\vartheta\kappa'$
implies $Q_\kappa(X,X)\le Q_{\kappa'}(X,X)$ for every $X$.

\emph{The bilinear identity.} For real $g_v$ on $\mathcal X$ put
$F(x)=\prod_{u\in V_+}g_u(x_u)$ and $G(y)=\prod_{u\in V_+}g_{\vartheta u}(y_u)$.
Since $V=V_+\sqcup\vartheta V_+$,
\begin{equation}\label{eq:cc:bilinear}
 \Lambda_\nu\big((g_v)_{v\in V}\big)=\sum_{x,y}F(x)\,C_\vartheta(\nu)_{x,y}\,G(y)
 =\ip{F}{C_\vartheta(\nu)G}=Q_\nu(F,G).
\end{equation}
When every $g_v$ equals $g$, both $F$ and $G$ are the half-function
$F_g(x)=\prod_{u\in V_+}g(x_u)$ carrying $g$ at every slot of the plus half; so
$\Lambda_\nu(g)=Q_\nu(F_g,F_g)$, at every cut of $V$.

\emph{Products.} If $\nu_1$, $\nu_2$ are finite signed measures on
$\mathcal X_1^V$, $\mathcal X_2^V$, we write $\nu_1\otimes\nu_2$ for the image of
the product measure under $(z^1,z^2)\mapsto((z^1_v,z^2_v))_{v\in V}$, a measure
on $(\mathcal X_1\times\mathcal X_2)^V$. Under the
identification $(\mathcal X_1\times\mathcal X_2)^{V_+}=\mathcal X_1^{V_+}\times
\mathcal X_2^{V_+}$,
\begin{equation}\label{eq:cc:kron}
 C_\vartheta(\nu_1\otimes\nu_2)=C_\vartheta(\nu_1)\otimes C_\vartheta(\nu_2)
 \qquad\text{(Kronecker product)},
\end{equation}
because the entry at $((x^1,x^2),(y^1,y^2))$ factors. A Kronecker product of
symmetric positive semidefinite matrices is symmetric positive semidefinite
(its eigenvalues are the products of those of the factors). Hence
$\kappa_i\succeq_\vartheta0$ for all $i$ implies
$\bigotimes_i\kappa_i\succeq_\vartheta0$.

\begin{remark}[Reflections, halves and permutations]\label{rem:cut-symmetries}
Let $\kappa$ be a finite signed measure on $\mathcal X^V$ and $\vartheta$ a cut.
\begin{enumerate}[label=(\alph*)]
\item \emph{Reflection.} $C_\vartheta(\vartheta_*\kappa)=C_\vartheta(\kappa)^{\mathsf T}$,
 since $\vartheta_*\kappa(z_{V_+}=x,\ z_{\vartheta V_+}=y)
 =\kappa(z_{\vartheta V_+}=x,\ z_{V_+}=y)$. As $\kappa$ is determined by its cut
 matrix, $\kappa$ is invariant under the reflection if and only if
 $C_\vartheta(\kappa)$ is symmetric. The \emph{symmetric} and \emph{skew parts}
 $\frac12(\kappa\pm\vartheta_*\kappa)$ have the symmetric and the antisymmetric
 part of $C_\vartheta(\kappa)$ as cut matrices. The diagonal cut form of a
 measure with antisymmetric cut matrix vanishes, and so does that of its
 product with any measure with symmetric cut matrix \eqref{eq:cc:kron}.
\item \emph{Halves.} $(\vartheta,V_-)$ is again a cut. Identify
 $\mathcal X^{V_-}$ with $\mathcal X^{V_+}$ by $x\mapsto(x_{\vartheta u})_{u\in V_+}$;
 then the cut matrix at $(\vartheta,V_-)$ is the transpose of the cut matrix
 at $(\vartheta,V_+)$. Hence $\kappa\succeq\kappa'$ at one choice of the plus
 half if and only if at the other; in particular reflection positivity does
 not depend on the choice of the plus half.
\item \emph{Permutations.} For a permutation $\sigma$ of $V$, the involution
 $\sigma^{-1}\vartheta\sigma$ with plus half $\sigma^{-1}V_+$ is a cut, and
 $s(x)=(x_{\sigma u})_{u\in\sigma^{-1}V_+}$ is a bijection
 $\mathcal X^{V_+}\to\mathcal X^{\sigma^{-1}V_+}$ with
 \[
  C_{\sigma^{-1}\vartheta\sigma}(\sigma_*\kappa)_{s(x),s(y)}=C_\vartheta(\kappa)_{x,y}.
 \]
 Indeed, for $w=z\circ\sigma$ we have $w_{\sigma^{-1}V_+}=s(z_{V_+})$ and
 $(w_{\sigma^{-1}\vartheta\sigma u})_{u\in\sigma^{-1}V_+}=s((z_{\vartheta v})_{v\in V_+})$.
 Hence $\kappa\succeq_\vartheta\kappa'$ if and only if
 $\sigma_*\kappa\succeq_{\sigma^{-1}\vartheta\sigma}\sigma_*\kappa'$.
\end{enumerate}
\end{remark}

\begin{lemma}[Common-cut order]\label{lem:commoncut}
Let $\vartheta$ be a cut of $V$ with halves $V_\pm$, and let
$\operatorname{Fix}(\vartheta)=\{z\in\mathcal X^V:z_{\vartheta v}=z_v\ \text{for
all }v\in V\}$.
\begin{enumerate}[label=(\roman*)]
\item \emph{Diagonal measures.} Let $\kappa$ be a nonnegative measure on
 $\mathcal X^V$ supported on $\operatorname{Fix}(\vartheta)$; we say that
 $\kappa$ is \emph{diagonal at $\vartheta$}, or \emph{fixed samplewise by
 $\vartheta$}. Then $C_\vartheta(\kappa)$ is diagonal with entries
 $\kappa(z_{V_+}=x)\ge0$; in particular $\kappa\succeq_\vartheta0$, and $\kappa$
 is RP at $\vartheta$. Every nonnegative measure whose support lies in that of
 $\kappa$ is again diagonal at $\vartheta$: a restriction of $\kappa$ to an
 event, a nonnegative reweighting of $\kappa$, and, if $\kappa=Z_*\pi$ for a
 finite set $\Omega$ with weights $\pi\ge0$ and a map
 $Z\colon\Omega\to\mathcal X^V$ (a \emph{latent representation}), every
 $Z_*(w\pi)$ with $w\ge0$ on $\Omega$ (a restriction to latent events if $w$ is
 an indicator). For a finite nonnegative measure $m$ on $\mathcal X^{V_+}$,
 the \emph{diagonal copy} $\operatorname{Diag}_\vartheta(m)$, the image of $m$
 under the map $x\mapsto z$ with $z_{V_+}=x$ and $z_{\vartheta u}=x_u$ for
 $u\in V_+$ (for a probability measure $m$: the law of $z$ with $z_{V_+}\sim m$
 and $z_{\vartheta u}=z_u$), is diagonal at $\vartheta$ with cut matrix
 $\operatorname{diag}(m)$.
\item \emph{Order.} If $\kappa,\kappa_1,\dots,\kappa_k$ are diagonal at
 $\vartheta$ and $\kappa\le\sum_i\kappa_i$, then
 $0\preceq_\vartheta\kappa\preceq_\vartheta\sum_i\kappa_i$.
\item \emph{Adjoining a factor.} If $\kappa\preceq_\vartheta\kappa'$ on
 $\mathcal X^V$ and $\rho$ is a finite signed measure on $\mathcal Y^V$ with
 $\rho\succeq_\vartheta0$, then $\kappa\otimes\rho\preceq_\vartheta\kappa'\otimes\rho$
 on $(\mathcal X\times\mathcal Y)^V$. In particular (ii) persists after
 adjoining $\rho$.
\item \emph{Diagonal functionals.} Consequently, in the situation of (ii) and
 (iii), $0\le\Lambda_{\kappa\otimes\rho}(g)\le\sum_i\Lambda_{\kappa_i\otimes\rho}(g)$
 for every real function $g$ on $\mathcal X\times\mathcal Y$.
\item \emph{One kernel on both halves.} Let $\vartheta'$ be a cut of a finite
 set $W$ with halves $W_\pm$, $\mathcal Y$ a finite set, and $P$ a Markov
 kernel from $\mathcal X^{V_+}$ to $\mathcal Y^{W_+}$ (a nonnegative matrix with
 unit row sums). For a finite signed measure $\kappa$ on $\mathcal X^V$ let
 $\kappa^P$ be the measure on $\mathcal Y^W$ obtained by applying $P$
 independently to $z_{V_+}$ and to the reflected minus half
 $(z_{\vartheta u})_{u\in V_+}$, and placing the results on $W_+$ and, through
 $\vartheta'$, on $W_-$:
 \[
  \kappa^P(w)=\sum_{z}\kappa(z)\,P\big(z_{V_+},w_{W_+}\big)\,
  P\big((z_{\vartheta u})_{u\in V_+},(w_{\vartheta'u'})_{u'\in W_+}\big).
 \]
 Then $C_{\vartheta'}(\kappa^P)=P^{\mathsf T}C_\vartheta(\kappa)P$; so
 $\kappa\succeq_\vartheta\kappa'$ implies $\kappa^P\succeq_{\vartheta'}\kappa'^P$.
 The measure $\kappa^P$ is nonnegative if $\kappa$ is, has the mass of
 $\kappa$, and has plus-half marginal $mP$ if $\kappa$ has plus-half marginal
 $m$. For a finite nonnegative measure $m$ on $\mathcal X^{V_+}$,
 $P^{\mathsf T}\operatorname{diag}(m)P\preceq\operatorname{diag}(mP)$; hence
 $\pi+k\operatorname{Diag}_\vartheta(m)\succeq_\vartheta0$ with $k\ge0$ implies
 $\pi^P+k\operatorname{Diag}_{\vartheta'}(mP)\succeq_{\vartheta'}0$. If $P$ is
 deterministic, $P(x,\cdot)=\delta_{\phi(x)}$, then $\kappa^P=\Phi_*\kappa$ for
 the \emph{half-local} map $\Phi\colon\mathcal X^V\to\mathcal Y^W$ with
 \[
  \Phi(z)_{W_+}=\phi(z_{V_+}),\qquad
  \big(\Phi(z)_{\vartheta'u'}\big)_{u'\in W_+}=\phi\big((z_{\vartheta u})_{u\in V_+}\big);
 \]
 the map $\Phi$ intertwines the two reflections, it maps
 $\operatorname{Fix}(\vartheta)$ into $\operatorname{Fix}(\vartheta')$, and
 $\Phi_*\operatorname{Diag}_\vartheta(m)=\operatorname{Diag}_{\vartheta'}(\phi_*m)$.
\end{enumerate}
\end{lemma}

\begin{proof}
(i) $C_\vartheta(\kappa)_{x,y}=\kappa(z_{V_+}=x,\ z_{\vartheta V_+}=y)$. On
$\operatorname{Fix}(\vartheta)$ we have $(z_{\vartheta u})_{u\in V_+}=z_{V_+}$, so
the entry vanishes unless $y=x$, and the diagonal entries are
$\kappa(z_{V_+}=x)\ge0$. A nonnegative diagonal matrix is symmetric and
positive semidefinite. The support of $Z_*(w\pi)$ is contained in
$Z(\{\pi>0\})$, which is the support of $\kappa$. The diagonal copy is
supported on $\operatorname{Fix}(\vartheta)$, with diagonal entries $m(x)$.

(ii) $\sum_i\kappa_i-\kappa$ is a nonnegative measure supported in
$\bigcup_i\operatorname{supp}\kappa_i\subseteq\operatorname{Fix}(\vartheta)$. By
linearity of $C_\vartheta$ and (i), $C_\vartheta(\sum_i\kappa_i)-C_\vartheta(\kappa)$
is nonnegative diagonal, and $C_\vartheta(\kappa)\succeq0$ by (i).

(iii) By \eqref{eq:cc:kron}, $C_\vartheta(\kappa'\otimes\rho)-C_\vartheta(\kappa\otimes\rho)
=C_\vartheta(\kappa'-\kappa)\otimes C_\vartheta(\rho)$, a Kronecker product of two
symmetric positive semidefinite matrices.

(iv) Apply \eqref{eq:cc:bilinear} with the same real vector $F_g$ on both sides
of the inequalities of (ii) and (iii).

(v) For $x',y'\in\mathcal Y^{W_+}$,
\[
 C_{\vartheta'}(\kappa^P)_{x',y'}
 =\sum_z\kappa(z)P\big(z_{V_+},x'\big)P\big((z_{\vartheta u})_u,y'\big)
 =\sum_{x,y}C_\vartheta(\kappa)_{x,y}P(x,x')P(y,y'),
\]
which is the matrix identity. Thus
$\ip{G}{C_{\vartheta'}(\kappa^P)G}=\ip{PG}{C_\vartheta(\kappa)PG}$ for every real
$G$, and $P^{\mathsf T}CP$ is symmetric when $C$ is; this gives the order
statement. Nonnegativity, the mass and the marginal follow from the unit row
sums. For the diagonal, by Jensen's inequality for each row of $P$,
\[
 \ip{G}{P^{\mathsf T}\operatorname{diag}(m)PG}=\sum_xm(x)\,(PG)(x)^2
 \le\sum_xm(x)\,(PG^2)(x)=\ip{G}{\operatorname{diag}(mP)G},
\]
so $C_{\vartheta'}(\pi^P)+k\operatorname{diag}(mP)\succeq
P^{\mathsf T}\big(C_\vartheta(\pi)+k\operatorname{diag}(m)\big)P\succeq0$, and
the left side is symmetric. If $P$ is deterministic, the defining sum
collapses to $\Phi_*\kappa$. Since $(z\circ\vartheta)_{V_+}=(z_{\vartheta u})_u$
and $((z\circ\vartheta)_{\vartheta u})_u=z_{V_+}$, we get
$\Phi(z\circ\vartheta)=\Phi(z)\circ\vartheta'$; for
$z\in\operatorname{Fix}(\vartheta)$ both halves of $\Phi(z)$ are $\phi(z_{V_+})$,
so $\Phi(z)\in\operatorname{Fix}(\vartheta')$. Finally
$P^{\mathsf T}\operatorname{diag}(m)P=\operatorname{diag}(\phi_*m)$, and two
measures with the same cut matrix coincide.
\end{proof}

\begin{remark}[Limits of the common-cut order]\label{rem:cc:limits}
\begin{enumerate}[label=(\alph*)]
\item Part (iv) is a union bound in the cone of positive semidefinite cut
 matrices, not a pointwise inequality for the integrand: a product of Fourier
 lifts can be negative. A Boolean cover $\one_B\le\sum_F\one_{I_F}$ of events of
 a latent representation therefore becomes the inequality
 $0\le\Lambda_{Z_*(\one_B\pi)\otimes\rho}(g)\le\sum_F\Lambda_{Z_*(\one_{I_F}\pi)\otimes\rho}(g)$
 of diagonal functionals.
\item The samplewise hypothesis is essential. Being supported on
 $\operatorname{Fix}(\vartheta)$ is a property of the samples; invariance of a
 measure under the reflection does not imply it, and reflection positivity of
 a measure at every cut does not make its restrictions to events positive
 semidefinite at a chosen cut without an additional argument. This is one
 reason why the residues modulo the small modulus and the position of an
 integer among the intervals are external indices of the comparison
 (\cref{def:hyb:coupled}) and not coordinates of the labels: a reference law
 carrying correlated residues would have to be restricted to events of the
 residues, and such restrictions need not remain positive semidefinite at the
 cuts where positivity is required.
\end{enumerate}
\end{remark}

\subsection{Controlling a law that is not reflection positive by two reflection-positive laws ($2\times2$ blocks)}\label{subsec:pos:brackets}

If a block matrix of two positive semidefinite diagonal blocks is positive
semidefinite, its off-diagonal block is bounded by the diagonal ones. Such a
\emph{bracket} controls a measure that is not reflection positive at a cut by
two positive ones, and brackets tensorize. This is how the comparison law,
which is not reflection positive, is controlled at the reserved cut.

\begin{lemma}[Block-Gram calculus]\label{lem:record:block-gram}
Let $\mathcal H,\mathcal H',\mathcal H_1,\dots,\mathcal H_n$ be finite-dimensional
Hilbert spaces.
\begin{enumerate}[label=(\alph*)]
\item \emph{Brackets.} For operators $M$ on $\mathcal H$, $M'$ on $\mathcal H'$
 and $K\colon\mathcal H'\to\mathcal H$, the block operator
 $\bigl(\begin{smallmatrix}M&K\\ K^*&M'\end{smallmatrix}\bigr)$ on
 $\mathcal H\oplus\mathcal H'$ is positive semidefinite if and only if
 $M\succeq0$, $M'\succeq0$ and $K=M^{1/2}CM'^{1/2}$ for an operator $C$ with
 $\|C\|\le1$. Then $|\ip{x}{Ky}|\le\ip{x}{Mx}^{1/2}\ip{y}{M'y}^{1/2}$.
\item \emph{Tensorization.} If
 $\bigl(\begin{smallmatrix}M_i&K_i\\ K_i^*&M_i'\end{smallmatrix}\bigr)\succeq0$
 on $\mathcal H_i\oplus\mathcal H_i$ for each $i$, then the same holds for
 $M=\bigotimes_iM_i$, $M'=\bigotimes_iM_i'$ and $K=\bigotimes_iK_i$, and
 $|\ip{x}{Ky}|\le\ip{x}{Mx}^{1/2}\ip{y}{M'y}^{1/2}$.
\item \emph{Modulus.} Every operator $T$ on a finite-dimensional Hilbert space
 satisfies
 \[
  \begin{pmatrix}|T^*|&T\\ T^*&|T|\end{pmatrix}\succeq0,
  \qquad |T|=(T^*T)^{1/2},\quad |T^*|=(TT^*)^{1/2}.
 \]
\item \emph{Symmetric brackets.} For self-adjoint $M$, $K$,
 $\bigl(\begin{smallmatrix}M&K\\ K&M\end{smallmatrix}\bigr)\succeq0$ if and only
 if $M\pm K\succeq0$. If $A_i$, $B_i$ are self-adjoint with $A_i\pm B_i\succeq0$
 for every $i$, then $\bigotimes_iA_i\pm\bigotimes_iB_i\succeq0$; for two
 factors, explicitly,
 \[
  2\,(A\otimes N\pm D\otimes G)=(A\pm D)\otimes(N+G)+(A\mp D)\otimes(N-G).
 \]
 In particular $0\preceq A_i\preceq A_i'$ for every $i$ implies
 $\bigotimes_iA_i\preceq\bigotimes_iA_i'$.
\end{enumerate}
\end{lemma}

\begin{proof}
(a) The diagonal blocks of a positive semidefinite block operator are positive
semidefinite. If $K=M^{1/2}CM'^{1/2}$ with $\|C\|\le1$, the form of the block
operator at $(x,y)$ is at least
$\|M^{1/2}x\|^2-2\|M^{1/2}x\|\,\|M'^{1/2}y\|+\|M'^{1/2}y\|^2\ge0$. Conversely,
for $\epsilon>0$ conjugate the block operator plus $\epsilon$ by
$\operatorname{diag}((M+\epsilon)^{-1/2},(M'+\epsilon)^{-1/2})$: the result
$\bigl(\begin{smallmatrix}I&C_\epsilon\\ C_\epsilon^*&I\end{smallmatrix}\bigr)$,
$C_\epsilon=(M+\epsilon)^{-1/2}K(M'+\epsilon)^{-1/2}$, is positive definite,
and its form at $(x,-C_\epsilon^*x)$ is $\|x\|^2-\|C_\epsilon^*x\|^2$, so
$\|C_\epsilon\|\le1$. Since $K=(M+\epsilon)^{1/2}C_\epsilon(M'+\epsilon)^{1/2}$
for every $\epsilon$, any limit point $C$ of $C_\epsilon$ as $\epsilon\to0$ is a
contraction with $K=M^{1/2}CM'^{1/2}$, by continuity of the square root. The
inequality is Cauchy--Schwarz: $|\ip{x}{Ky}|=|\ip{M^{1/2}x}{CM'^{1/2}y}|$.

(b) Write $K_i=M_i^{1/2}C_iM_i'^{1/2}$ by (a). The square root of a tensor
product of positive semidefinite operators is the tensor product of the square
roots, so $K=M^{1/2}\big(\bigotimes_iC_i\big)M'^{1/2}$, and
$\|\bigotimes_iC_i\|=\prod_i\|C_i\|\le1$; apply (a).

(c) Let $T=U|T|$ with $U$ unitary. Then $|T^*|=U|T|U^*$, and the block operator
equals $X|T|X^*$, where $X$ is the column operator $h\mapsto(Uh,h)$ into
$\mathcal H\oplus\mathcal H$; indeed $X^*(a,b)=U^*a+b$, $U|T|=T$ and
$|T|U^*=T^*$.

(d) Conjugation by the unitary
$\frac1{\sqrt2}\bigl(\begin{smallmatrix}I&I\\ I&-I\end{smallmatrix}\bigr)$ turns
$\bigl(\begin{smallmatrix}M&K\\ K&M\end{smallmatrix}\bigr)$ into
$\bigl(\begin{smallmatrix}M+K&0\\ 0&M-K\end{smallmatrix}\bigr)$. If
$A_i\pm B_i\succeq0$, the first part gives the brackets
$\bigl(\begin{smallmatrix}A_i&B_i\\ B_i&A_i\end{smallmatrix}\bigr)\succeq0$, (b)
tensorizes them, $\bigotimes_iB_i$ is self-adjoint, and the first part applies
again. The two-factor identity is expanded directly, and each summand on its
right is a tensor product of positive semidefinite operators. Finally
$0\preceq A_i\preceq A_i'$ gives $A_i'\pm A_i\succeq0$.
\end{proof}

\emph{Brackets of measures.} For finite signed measures $\kappa,\kappa',\zeta$
on $\mathcal X^V$ and a cut $\vartheta$ we write
$\bigl(\begin{smallmatrix}\kappa&\zeta\\ \zeta^*&\kappa'\end{smallmatrix}\bigr)\succeq_\vartheta0$,
a \emph{bracket at $\vartheta$}, if the block matrix with blocks
$C_\vartheta(\kappa)$, $C_\vartheta(\zeta)$, $C_\vartheta(\zeta)^{\mathsf T}$ (the block denoted $\zeta^*$; $\zeta^*$ is not a separate measure),
$C_\vartheta(\kappa')$ is positive semidefinite. By (a),
$|Q_\zeta(X,Y)|\le Q_\kappa(X,X)^{1/2}Q_{\kappa'}(Y,Y)^{1/2}$ for all
half-functions; by (b) and \eqref{eq:cc:kron}, brackets of the factors of
product measures tensorize to a bracket of the products; and by (d), if
$C_\vartheta(\zeta)$ is symmetric, then
$\bigl(\begin{smallmatrix}\kappa&\zeta\\ \zeta&\kappa\end{smallmatrix}\bigr)\succeq_\vartheta0$
if and only if $\kappa\pm\zeta\succeq_\vartheta0$.

\subsection{From one cut to all slots: the H\"older inequality}\label{subsec:pos:holder}

Reflection positivity turns a quadratic-form inequality at one cut into a
mixed inequality across all slots. We first prove the single-cut
Cauchy--Schwarz step, then explain why a sorting sequence gives the signed
H\"older inequality. The method comes from infrared and chessboard estimates in
statistical mechanics
\cite{FrohlichSimonSpencer76,DysonLiebSimon78,FrohlichLieb,FILS,GlimmJaffe87};
the same positive semidefinite formalism appears in graph-density inequalities
and flag algebras
\cite{ConlonLee,HatamiNorine11,Razborov07,LovaszGraphLimits12,BorgsChayesLovaszSosVesztergombi08}.
We use only the finite cut-matrix form proved below.

\begin{definition}[Folds; sorting sequences]\label{def:sorting}
For a cut $\vartheta$ define $\phi^\vartheta_\pm\colon V\to V$ by
$\phi^\vartheta_+(v)=v$ for $v\in V_+$ and $\phi^\vartheta_+(v)=\vartheta v$ for
$v\in V_-$; and $\phi^\vartheta_-(v)=\vartheta v$ for $v\in V_+$ and
$\phi^\vartheta_-(v)=v$ for $v\in V_-$. A finite sequence of cuts
$\vartheta^1,\dots,\vartheta^T$ with chosen signs $s_t\in\{\pm\}$ is
\emph{sorting} if
$\phi^{\vartheta^T}_{s_T}\circ\dots\circ\phi^{\vartheta^1}_{s_1}$ is constant on
$V$.
\end{definition}

\begin{lemma}[Cauchy--Schwarz at one cut]\label{lem:holder:one-cut}
Suppose $\nu$ is reflection positive at $\vartheta$. Then for every assignment
$h\colon V\to\R^{\mathcal X}$,
\begin{equation}\label{eq:cs-cut}
 |\Lambda_\nu(h)|^2\le\Lambda_\nu(h\circ\phi^\vartheta_+)\,
 \Lambda_\nu(h\circ\phi^\vartheta_-),\qquad
 \Lambda_\nu(h\circ\phi^\vartheta_\pm)\ge0 .
\end{equation}
In particular $\Lambda_\nu(g)\ge0$ for a constant assignment.
\end{lemma}

\begin{proof}
Put $F(x)=\prod_{u\in V_+}h(u)(x_u)$ and $G(x)=\prod_{u\in V_+}h(\vartheta u)(x_u)$
for $x\in\mathcal X^{V_+}$. By \eqref{eq:cc:bilinear}, $\Lambda_\nu(h)$ is the
cut form of $\nu$ at $(F,G)$, and its values at $(F,F)$ and $(G,G)$ are
$\Lambda_\nu(h\circ\phi^\vartheta_+)$ and $\Lambda_\nu(h\circ\phi^\vartheta_-)$:
the first folded assignment puts $h(u)$ at $u$ and at $\vartheta u$, the second
puts $h(\vartheta u)$ at both ($u\in V_+$). Cauchy--Schwarz for the positive
semidefinite form of $C_\vartheta(\nu)$ gives \eqref{eq:cs-cut}. For a constant
assignment $F=G$.
\end{proof}

\begin{theorem}[Signed H\"older inequality]\label{thm:holder}
Let $\nu$ be a finite signed measure on $\mathcal X^V$ that is RP at every cut
of a sorting sequence. Then for all real functions $g$, $g_v$ on $\mathcal X$,
\[
 \Lambda_\nu(g)\ge0\qquad\text{and}\qquad
 \big|\Lambda_\nu\big((g_v)_{v\in V}\big)\big|\le\prod_{v\in V}\Lambda_\nu(g_v)^{1/d}.
\]
\end{theorem}

\begin{proof}
By \cref{lem:holder:one-cut}, diagonal values are nonnegative. Fix real
functions $f_1,\dots,f_k$ and let $M=\max_h|\Lambda_\nu(h)|$ over the
assignments $h\colon V\to\{f_1,\dots,f_k\}$. If $M>0$, let $h$ attain it. Both
folded assignments again take values in $\{f_i\}$, so by \eqref{eq:cs-cut}
they are nonnegative, at most $M$, and have product at least $M^2$. Hence both
equal $M$: folds of maximisers are maximisers. Applying the folds in the order
$\vartheta^T,\vartheta^{T-1},\dots,\vartheta^1$ turns $h$ into
$h\circ\phi^{\vartheta^T}_{s_T}\circ\dots\circ\phi^{\vartheta^1}_{s_1}$, which is a
constant maximiser. Therefore $M=\max_i\Lambda_\nu(f_i)$, which also holds when
$M=0$. For $\epsilon>0$ apply this to
$f_v=g_v/(\Lambda_\nu(g_v)+\epsilon)^{1/d}$, which satisfy $\Lambda_\nu(f_v)\le1$.
Multilinearity gives
$|\Lambda_\nu((g_v))|\le\prod_v(\Lambda_\nu(g_v)+\epsilon)^{1/d}$; let
$\epsilon\downarrow0$.
\end{proof}

The proof uses only reflection positivity at the cuts of the sequence; it does
not use that $\nu$ is nonnegative or normalized. Recall that a
\emph{seminorm} is a nonnegative, absolutely homogeneous function satisfying
the triangle inequality; it may vanish on nonzero vectors.

\begin{corollary}[Seminorm]\label{cor:seminorm}
Under the hypotheses of \cref{thm:holder}, $\|g\|_\nu=\Lambda_\nu(g)^{1/d}$ is a
seminorm on $\R^{\mathcal X}$. If every one-slot marginal of $\nu$ is the
uniform probability measure on $\mathcal X$, then
$|\E_{x\in\mathcal X}g(x)|\le\|g\|_\nu$.
\end{corollary}

\begin{proof}
Homogeneity holds because $d$ is even. For the triangle inequality, expand
$\Lambda_\nu(g_1+g_2)$ into $2^d$ mixed terms and bound each by
\cref{thm:holder}:
\[
 0\le\Lambda_\nu(g_1+g_2)\le\sum_{S\subseteq V}\|g_1\|_\nu^{|S|}\|g_2\|_\nu^{d-|S|}
 =(\|g_1\|_\nu+\|g_2\|_\nu)^d .
\]
For the mean, put $g$ in one slot and $1$ in the others; then $\Lambda_\nu(1)=1$
and the uniform marginal give
$|\E g|=|\Lambda_\nu(g,1,\dots,1)|\le\|g\|_\nu\|1\|_\nu^{d-1}=\|g\|_\nu$.
\end{proof}

\begin{remark}[Domination at one cut]\label{rem:cc:mixed}
Let $\mu$ satisfy the hypotheses of \cref{thm:holder}, let $\vartheta$ be a cut
at which $\mu$ is RP (it need not belong to the sorting sequence), and let $D$
be a finite signed measure with a bracket
$\bigl(\begin{smallmatrix}\mu&D\\ D^*&\mu\end{smallmatrix}\bigr)\succeq_\vartheta0$,
for instance $0\preceq_\vartheta D\preceq_\vartheta\mu$
(\cref{lem:record:block-gram}(d)). With $F,G$ as in \eqref{eq:cc:bilinear},
\cref{lem:record:block-gram}(a) and \cref{thm:holder} for the two folded
assignments, in which each $g_u$ with $u\in V_+$, respectively $u\in V_-$,
occurs twice, give
\[
 |\Lambda_D((g_v)_v)|\le Q_\mu(F,F)^{1/2}Q_\mu(G,G)^{1/2}
 \le\prod_{v\in V}\Lambda_\mu(g_v)^{1/d}.
\]
The measure $D$ need not be positive, RP at any other cut, a product, or
normalized, and no seminorm is attached to it; only the dominating law
carries the H\"older inequality. Normalizing $D$ by its mass would in general
destroy the domination, and we never do it.
\end{remark}

\subsection{The $72$ slots $(\Z/6\Z)^2\times\{0,1\}$, their symmetries, and the six queries}\label{subsec:slot-array}

The slots of the functional are indexed by a product of reflection
coordinates: two six-chamber coordinates and a two-point layer. A product of
coordinates, each with its own cuts and its own sorting sequence, is sorted by
the concatenation of the sequences, because folds of different coordinates
commute; this is the reason for working with products. A chamber coordinate
carries the rotation $\rot$, the relation $\Rel$ and the group $\mathcal G_6$ of chamber automorphisms defined below, which transport
positivity from one chamber cut to the others.

\begin{definition}[Reflection coordinates; the slot array]\label{def:slot-array}
\leavevmode
\begin{enumerate}[label=(\alph*)]
\item A \emph{reflection coordinate} is a finite set with a family of cuts that
 contains a sorting sequence. For reflection coordinates $V^{(i)}$,
 $i\in\mathcal C$ (a finite index set), a cut $\vartheta$ of $V^{(i)}$ with plus
 half $V^{(i)}_+$ induces the cut $\vartheta^{[i]}$ of $V=\prod_{i\in\mathcal C}V^{(i)}$
 that applies $\vartheta$ to the coordinate $v_i$ and fixes the other
 coordinates, with plus half $\{v:v_i\in V^{(i)}_+\}$; these are the \emph{cuts
 of the coordinate $i$}.
\item The \emph{chamber coordinate} is $\Zsix$, whose elements are
 \emph{chambers}, with the \emph{rotation} $\rot(j)=j+2$ for even $j$ and
 $\rot(j)=j-2$ for odd $j$, the relation
 $\Rel(x,y)\Leftrightarrow y\in\{x,x+1,x-1,\rot(x)\}$, and, for $a\in\Zsix$,
 the \emph{chamber reflection} $\vartheta_a(j)=2a-1-j$ with plus half
 $H_a=\{a,a+1,a+2\}$; its family of cuts consists of the six \emph{chamber
 cuts} $(\vartheta_a,H_a)$, $a\in\Zsix$ (they are cuts by
 \cref{lem:slot-array}(i) and contain a sorting sequence by
 \cref{lem:slot-array}(ii)). A pair $(j,\rot j)$ is a \emph{$\rot$-step}. The
 plus fold of $\vartheta_a$ is $f_a$: $f_a(j)=j$ for $j\in H_a$ and
 $f_a(j)=\vartheta_a(j)$ otherwise.
\item The \emph{layer} is $\{0,1\}$ with the single cut $\ell\mapsto1-\ell$ and
 plus half $\{0\}$; its plus fold $\ell\mapsto0$ is constant, so this cut is a
 sorting sequence.
\item The \emph{slot array} is $V=\Zsix\times\Zsix\times\{0,1\}$, $d=|V|=72$,
 with slots $v=(u_0,u_1,\ell)$; the coordinates $0$ and $1$ are the
 \emph{chamber coordinates} and the coordinate $2$, with value $\ell$, is the
 \emph{layer}. The \emph{chamber cut} $\vartheta_{c,a}$ ($c\in\{0,1\}$,
 $a\in\Zsix$) is the cut of the coordinate $c$ induced by $\vartheta_a$. The
 \emph{reserved cut} is the cut of the layer,
 \[
  \iota(u_0,u_1,\ell)=(u_0,u_1,1-\ell),\qquad V_+=\{\ell=0\},\qquad V_-=\iota V_+ .
 \]
 The \emph{family of cuts} of $V$ is
 $\mathfrak F=\{\vartheta_{c,a}:c\in\{0,1\},\ a\in\Zsix\}\cup\{\iota\}$: thirteen cuts, with seven distinct
 involutions, since $\vartheta_{a+3}=\vartheta_a$ with complementary halves
 (three involutions in each chamber coordinate, and $\iota$). At
 $\iota$ we write $C_\iota$, $\succeq_\iota$ and $Q_\kappa$ for the cut matrix,
 the cut order and the cut form (\cref{def:cc:order}).
\item The \emph{queries} are the six ordered pairs of slots
 \[
  e^{(j)}=\big((0,j,0),\,(\rot0,\rot j,1)\big)=\big((0,j,0),\,(2,\rot j,1)\big),
  \qquad j\in\Zsix,
 \]
 forming the family $\mathcal E_6$. For $e=(u_e,v_e)\in\mathcal E_6$, $u_e$ is
 the \emph{source} and $v_e$ the \emph{target}.
\end{enumerate}
\end{definition}

\begin{lemma}[The slot array]\label{lem:slot-array}
\leavevmode
\begin{enumerate}[label=(\roman*)]
\item \emph{Chamber geometry and transport.} Each $\vartheta_a$ is a
 fixed-point-free involution of $\Zsix$ commuting with $\rot$, with
 $\vartheta_a(H_a)=\Zsix\setminus H_a$; moreover $\vartheta_{a+3}=\vartheta_a$
 and $H_{a+3}=\Zsix\setminus H_a$, so $\vartheta_{c,a+3}$ is $\vartheta_{c,a}$
 with its halves exchanged. The \emph{chamber automorphisms}, the
 permutations $\alpha$ of $\Zsix$ with $\alpha\circ\rot=\rot\circ\alpha$ and
 $\Rel(\alpha x,\alpha y)\Leftrightarrow\Rel(x,y)$, form a group $\mathcal G_6$
 of order $6$: the translations $j\mapsto j+2k$ ($k=0,1,2$) and the
 reflections $\vartheta_0,\vartheta_1,\vartheta_2$. For $a\in\Zsix$ the element
 $\alpha_a\in\mathcal G_6$ given by $\alpha_a(j)=j-a$ for even $a$ and
 $\alpha_a(j)=a+2-j$ for odd $a$ satisfies $\alpha_a(H_a)=H_0$ and
 $\alpha_a\vartheta_a\alpha_a^{-1}=\vartheta_0$. There is exactly one bijection
 $H_a\to H_0$ that preserves $\Rel$ (the \emph{half isomorphism} of $H_a$); it
 is the restriction of every $\alpha\in\mathcal G_6$ with $\alpha(H_a)=H_0$, and
 it maps the unique $\rot$-step inside $H_a$ to $(0,2)$. \emph{Transport:} let
 $V=\prod_{i\in\mathcal C}V^{(i)}$ be a product of reflection coordinates whose
 coordinate $c$ is a chamber coordinate (for instance $\Zsix$ itself,
 $\Zsix\times\{0,1\}$, or the slot array with $c\in\{0,1\}$), write
 $\vartheta_{c,a}$ for the cut of the coordinate $c$ induced by $\vartheta_a$, and
 for $\alpha\in\mathcal G_6$ let $\alpha^{(c)}$ be the permutation of $V$ that
 applies $\alpha$ to the coordinate $c$. If a finite signed measure $\kappa$ on
 $\mathcal X^V$ is invariant under $\alpha_a^{(c)}$, then
 $C_{\vartheta_{c,a}}(\kappa)=S_aC_{\vartheta_{c,0}}(\kappa)S_a^{\mathsf T}$
 for the permutation matrix $S_a$ of the bijection
 $x\mapsto(x_{\alpha_a^{(c)}u})_u$ of half-labellings; so a lower bound
 $C_{\vartheta_{c,0}}(\kappa)\succeq Q$ gives
 $C_{\vartheta_{c,a}}(\kappa)\succeq S_aQS_a^{\mathsf T}$, and
 $\kappa\succeq_{\vartheta_{c,0}}\kappa'$ implies
 $\kappa\succeq_{\vartheta_{c,a}}(\alpha_a^{(c)})_*\kappa'$. In particular a
 measure invariant under $\mathcal G_6$ acting on the coordinate $c$ that is RP
 at $\vartheta_{c,0}$ is RP at every chamber cut of the coordinate $c$; and on
 the slot array a measure invariant under $\mathcal G_6$ acting on each chamber
 coordinate that is RP at $\vartheta_{0,0}$, $\vartheta_{1,0}$ and $\iota$ is RP
 at every cut of $\mathfrak F$.
\item \emph{Folds.} $f_0=(0,1,2,2,1,0)$ and $f_4=(0,0,5,4,4,5)$ (values at
 $0,\dots,5$), and $f_0\circ f_4\circ f_0$ is the constant map $0$; so
 $(\vartheta_0,\vartheta_4,\vartheta_0)$, all with sign $+$, is a sorting
 sequence of $\Zsix$. The minus fold of $\vartheta_a$ is $f_{a+3}$; in
 particular $f_4$ is the plus fold of $\vartheta_4$, the involution
 $\vartheta_1=\vartheta_4$ with plus half $H_4=\{4,5,0\}$.
\item \emph{Products and sorting.} In a product $V=\prod_{i\in\mathcal C}V^{(i)}$
 of reflection coordinates every induced cut is a cut; its folds act on its
 own coordinate by the folds of the factor and fix the other coordinates;
 reflections and folds of cuts of different coordinates commute; and every
 interleaving of sorting sequences of the factors that keeps the order within
 each factor is a sorting sequence of $V$. For the slot array, $\iota$ is a
 cut, every chamber reflection and chamber fold commutes with $\iota$ and its
 folds, the seven-step sequence
 \[
  \vartheta_{0,0},\ \vartheta_{0,4},\ \vartheta_{0,0},\ \vartheta_{1,0},\
  \vartheta_{1,4},\ \vartheta_{1,0},\ \iota,
 \]
 all with sign $+$, maps every slot to $(0,0,0)$, and $\mathfrak F$ consists
 of seven distinct involutions. Consequently a finite signed measure $\nu$ on
 $\mathcal X^V$ that is RP at the cuts of this sequence satisfies
 \cref{thm:holder,cor:seminorm} with $d=72$; if its one-slot marginals are
 uniform probabilities, $(\E g)^{72}\le\Lambda_\nu(g)$ whenever $\E g\ge0$,
 which is $\alpha^{72}\le Y$ for the functional (\cref{lem:record:holder}(e)).
\item \emph{Queries.} The twelve endpoints of the six queries are distinct;
 every source lies in $V_+$ and every target in $V_-$; the reflected targets
 $\iota v_e=(2,\rot j,0)$ lie in $V_+$; and the twelve slots $u_e$, $\iota v_e$
 ($e\in\mathcal E_6$) of $V_+$ are distinct. Each query applies $\rot$ to both
 chamber coordinates and flips the layer.
\item \emph{Cylinder pullback.} Let $V=\prod_{i\in\mathcal C}V^{(i)}$ be a product
 of reflection coordinates, $\emptyset\ne\mathcal C'\subseteq\mathcal C$ a set
 of \emph{active} coordinates, $W=\prod_{i\in\mathcal C'}V^{(i)}$ with the
 induced cuts, and $\pi\colon V\to W$ the projection. The \emph{pullback}
 $\rho^*$ of a finite signed measure $\rho$ on $\mathcal X^W$ is its image under
 $\zeta\mapsto\zeta\circ\pi$: the slot $v$ receives the label of $\pi(v)$.
 \begin{enumerate}[label=(\alph*)]
 \item \emph{Inactive cuts.} If $\vartheta$ is a cut of a coordinate
  $i\notin\mathcal C'$, every $\zeta\circ\pi$ lies in
  $\operatorname{Fix}(\vartheta^{[i]})$; so every nonnegative measure supported
  on such labellings, for instance the pullback of any nonnegative measure on
  $\mathcal X^W$, is diagonal at $\vartheta^{[i]}$ and has a nonnegative
  diagonal cut matrix there (\cref{lem:commoncut}(i)).
 \item \emph{Active cuts.} If $\vartheta$ is a cut of a coordinate
  $i\in\mathcal C'$, then $\zeta\mapsto\zeta\circ\pi$ is half-local for the cuts
  induced by $\vartheta$ on $W$ and on $V$ (\cref{lem:commoncut}(v)); so an
  inequality $\rho_1\succeq\rho_2$ at the cut of $W$ gives $\rho_1^*\succeq\rho_2^*$
  at the cut of $V$, and in particular reflection positivity passes from
  $\rho$ to $\rho^*$.
 \item \emph{Marginals.} For $u,v\in V$ the law of $(z_u,z_v)$ under $\rho^*$ is
  the law of $(\zeta_{\pi u},\zeta_{\pi v})$ under $\rho$.
 \item \emph{Equivariance.} If $\rho$ is invariant under
  $\zeta\mapsto(g(\zeta_w))_w$ for a bijection $g$ of $\mathcal X$, then $\rho^*$
  is invariant under $z\mapsto(g(z_v))_v$.
 \end{enumerate}
 For the slot array and $\mathcal C'=\{0,2\}$ (a law on labels indexed by
 $(u_0,\ell)$), the cuts of the second chamber coordinate are inactive, and by
 (c) the ordered pair of labels at every query of $\mathcal E_6$ under the
 pullback $\rho^*$ has the law of the single ordered pair at the slots
 $(0,0)$ and $(2,1)$ of $\Zsix\times\{0,1\}$ under the original law $\rho$.
\end{enumerate}
\end{lemma}

\begin{proof}
(i) $\vartheta_a(\vartheta_a(j))=2a-1-(2a-1-j)=j$. A fixed point would satisfy
$2j\equiv2a-1\pmod6$, which is impossible because $2j-2a+1$ is odd. For
$i\in\{0,1,2\}$, $\vartheta_a(a+i)=a-1-i$, so
$\vartheta_a(H_a)=\{a-1,a-2,a-3\}=\Zsix\setminus H_a$; and
$\vartheta_{a+3}(j)=2a+5-j\equiv\vartheta_a(j)$,
$H_{a+3}=\{a+3,a+4,a+5\}=\Zsix\setminus H_a$. The map $\vartheta_a$ reverses
parity. For even $j$, $\vartheta_a(j)$ is odd and
$\rot(\vartheta_a(j))=2a-3-j=\vartheta_a(j+2)=\vartheta_a(\rot j)$; for odd $j$,
$\rot(\vartheta_a(j))=2a+1-j=\vartheta_a(j-2)=\vartheta_a(\rot j)$.

\emph{Chamber automorphisms.} Let $\alpha$ be one. For every $x$, the elements
$y\ne x$ with $\Rel(x,y)$ and $y\ne\rot(x)$ are $x\pm1$, because
$\rot(x)=x\pm2$. Since $\alpha$ is a bijection commuting with $\rot$ and
preserving $\Rel$, it maps this set for $x$ onto the set for $\alpha x$; so it
preserves the hexagon adjacency $y=x\pm1$, and $\alpha$ is a rotation
$j\mapsto j+k$ or a reflection $j\mapsto m-j$. A rotation commutes with $\rot$
if and only if $k$ is even: then it preserves parity and
$\rot(j+k)=\rot(j)+k$, while for odd $k$, $\rot(k)=k-2\ne k+2=\rot(0)+k$. A
reflection commutes with $\rot$ if and only if $m$ is odd: then it reverses
parity and $\rot(m-j)=m-\rot(j)$, while for even $m$,
$\rot(m)=m+2\ne m-2=m-\rot(0)$. The odd values $m=2a-1$ give
$\vartheta_0,\vartheta_1,\vartheta_2$. Conversely these six maps commute with
$\rot$ and change the differences $y-x\in\{0,\pm1\}$ at most by a sign, so
they preserve $\Rel$, and so do their inverses. The chamber automorphisms form
a group by definition; it has these six elements.

\emph{The elements $\alpha_a$.} For even $a$, $\alpha_a$ is a translation by an
even number, $\alpha_a(H_a)=H_0$, and
$\alpha_a\vartheta_a\alpha_a^{-1}(j)=2a-1-(j+a)-a=-1-j=\vartheta_0(j)$. For odd
$a$, $\alpha_a$ is the reflection $j\mapsto m-j$ with $m=a+2$ odd, hence lies
in $\mathcal G_6$; it is an involution with $\alpha_a(H_a)=\{2,1,0\}$, and
$\alpha_a\vartheta_a\alpha_a(j)=a+2-\big(2a-1-(a+2-j)\big)=5-j=\vartheta_0(j)$.

\emph{Half isomorphisms.} Restricted to $H_0=\{0,1,2\}$, the relation $\Rel$
holds on all ordered pairs except $(2,0)$: indeed $\rot(0)=2$, while
$0\notin\{2,1,3,4\}$. So $(0,2)$ is the unique pair $(x,y)$ of $H_0$ with
$\Rel(x,y)$ and not $\Rel(y,x)$, and it is the only $\rot$-step inside $H_0$
($\rot(1)=5$ and $\rot(2)=4$ lie outside $H_0$). The automorphism $\alpha_a$
maps $H_a$ onto $H_0$ preserving $\Rel$ and $\rot$-steps, so $H_a$ also has a
unique asymmetric pair, which is its unique $\rot$-step. A $\Rel$-preserving
bijection $H_a\to H_0$ must map the asymmetric pair of $H_a$ to $(0,2)$ and
therefore the third element to $1$. This proves uniqueness, and the restriction
of every $\alpha\in\mathcal G_6$ with $\alpha(H_a)=H_0$ is such a bijection.

\emph{Transport.} Let $\sigma=\alpha_a^{(c)}$. Since
$\alpha_a^{-1}\vartheta_0\alpha_a=\vartheta_a$, we have
$\sigma^{-1}\vartheta_{c,0}\sigma=\vartheta_{c,a}$, and
$\sigma^{-1}\{v:v_c\in H_0\}=\{v:v_c\in\alpha_a^{-1}H_0\}=\{v:v_c\in H_a\}$ is
the plus half of $\vartheta_{c,a}$. If $\sigma_*\kappa=\kappa$, then
\cref{rem:cut-symmetries}(c) gives the relabelling of the cut matrix, and
$\kappa\succeq_{\vartheta_{c,0}}\kappa'$ implies
$\kappa=\sigma_*\kappa\succeq_{\vartheta_{c,a}}\sigma_*\kappa'$. For the last
two assertions apply this with $\kappa'=0$ to each $a\in\Zsix$ and each chamber
coordinate.

(ii) For $f_0$: $H_0=\{0,1,2\}$ and $\vartheta_0(3,4,5)=(2,1,0)$. For $f_4$:
$H_4=\{4,5,0\}$ and $\vartheta_4(j)=7-j\equiv1-j$, so $\vartheta_4(1,2,3)=(0,5,4)$.
Hence
\[
 f_0(\Zsix)=\{0,1,2\},\qquad f_4(\{0,1,2\})=\{0,5\},\qquad f_0(\{0,5\})=\{0\},
\]
and $f_0\circ f_4\circ f_0$ is the constant map $0$. The minus fold of
$\vartheta_a$ maps $H_a$ by $\vartheta_a$ onto $H_{a+3}$ and fixes $H_{a+3}$; as
$\vartheta_{a+3}=\vartheta_a$, it is the plus fold $f_{a+3}$ of $\vartheta_{a+3}$.

(iii) An induced cut $\vartheta^{[i]}$ is a fixed-point-free involution because
$\vartheta$ is one, and it maps $\{v:v_i\in V^{(i)}_+\}$ onto
$\{v:v_i\in V^{(i)}_-\}$. If $v_i\in V^{(i)}_+$, then $v$ lies in the plus half
and is fixed by the plus fold; otherwise the plus fold sends $v$ to
$\vartheta^{[i]}v$, whose coordinate $i$ is $\vartheta(v_i)$. So the folds of
$\vartheta^{[i]}$ replace $v_i$ by its fold and fix the other coordinates, and
maps acting on different coordinates commute. In the composition of the folds
of an interleaved sequence, the coordinate $i$ is changed only by the folds of
the coordinate $i$, in their order, and their composition is constant on
$V^{(i)}$; so the composition is constant on $V$. For the slot array the factor
sequences are $(\vartheta_0,\vartheta_4,\vartheta_0)$ for both chamber
coordinates, by (ii), and the single cut of the layer, whose plus fold
$\ell\mapsto0$ is constant; $\iota$ is the cut of the layer. The induced
involutions $\vartheta_{c,a}$ and $\vartheta_{c,a'}$ coincide if and only if
$\vartheta_a=\vartheta_{a'}$, that is $2a\equiv2a'\pmod6$, that is
$a'\in\{a,a+3\}$; involutions of different coordinates differ; so $\mathfrak F$
consists of $3+3+1=7$ involutions. The last assertion is
\cref{thm:holder,cor:seminorm}.

(iv) The sources $(0,j,0)$ are six distinct slots with $u_0=0$ and $\ell=0$; the
targets $(2,\rot j,1)$ are six distinct slots, since $\rot$ is a bijection, with
$u_0=2$ and $\ell=1$. As $0\ne2$, the twelve endpoints are distinct. Sources lie
in $V_+$ and targets in $V_-$, and $\iota v_e=(2,\rot j,0)\in V_+$; the slots
$u_e$ and $\iota v_e$ are distinct, since their first coordinates are $0$ and
$2$. The last sentence is the definition, with $\rot0=2$.

(v) (a) If $i\notin\mathcal C'$, then $\vartheta^{[i]}v$ differs from $v$ only in
the coordinate $i$, so $\pi(\vartheta^{[i]}v)=\pi(v)$ and
$(\zeta\circ\pi)_{\vartheta^{[i]}v}=(\zeta\circ\pi)_v$. (b) Write $\vartheta_V$,
$\vartheta_W$ for the cuts induced by $\vartheta$ on $V$ and on $W$, with plus
halves $V_+$, $W_+$. Since $i\in\mathcal C'$, $\pi\circ\vartheta_V=\vartheta_W\circ\pi$,
$\pi(V_+)\subseteq W_+$ and $\pi^{-1}(W_+)=V_+$. Put
$\phi(\xi)=(\xi_{\pi v})_{v\in V_+}$ for $\xi\in\mathcal X^{W_+}$. For
$z=\zeta\circ\pi$ we have $z_{V_+}=\phi(\zeta_{W_+})$ and
$(z_{\vartheta_Vv})_{v\in V_+}=(\zeta_{\vartheta_W\pi v})_{v\in V_+}
=\phi\big((\zeta_{\vartheta_Ww})_{w\in W_+}\big)$; so the pullback is the
half-local map of $\phi$, and \cref{lem:commoncut}(v) applies. (c) and (d) are
immediate from $z_v=\zeta_{\pi v}$. For the slot array with $\mathcal C'=\{0,2\}$,
$\pi(0,j,0)=(0,0)$ and $\pi(2,\rot j,1)=(2,1)$.
\end{proof}

\begin{remark}\label{rem:slot-array-instances}
Nothing in the transport of (i), in (iii) and in (v) depends on the particular
factors. For instance the
array $(\Zsix)^3$, with the nine-step sequence formed from
$(\vartheta_0,\vartheta_4,\vartheta_0)$ in each of its three chamber
coordinates, is sorted in the same way; Part~VI uses this array.
\end{remark}

\subsection{Square-affine invariance}\label{subsec:pos:invariance}

Many laws below are invariant under relabelling all field elements
at once by maps $x\mapsto rx+c$ with $r$ a square unit, chosen
independently at each prime. Such a map acts only on the sampled
field labels, not on the discrete choices in a record (which
component or type was selected). Conditioning on a component or a
type therefore preserves this invariance, while conditioning on a
specified field label need not preserve it, and no invariance is claimed
after such conditioning below. This subsection
defines the invariance and lists the operations that preserve it.

\begin{definition}[Square-affine invariance and pair law]\label{def:map:invariance}
Let $\mathcal B$ be a finite set of odd primes with exponents $e_p\ge1$
($p\in\mathcal B$), and $X_{\mathcal B}=\prod_{p\in\mathcal B}\Z/p^{e_p}\Z$. Let
$\mathfrak G_{\mathcal B}$ be the group of maps $g=(g_p)_{p\in\mathcal B}$ of
$X_{\mathcal B}$, $g(x)_p=r_px_p+c_p$, where $r_p$ is a square unit modulo
$p^{e_p}$ and $c_p\in\Z/p^{e_p}\Z$, chosen \emph{independently} for each prime
(for $e_p=1$: $r_p\in\Qsq$ and $c_p\in\Fp$). It acts on tuples
$z\in X_{\mathcal B}^S$ ($S$ any finite index set) by $(gz)_s=g(z_s)$, the same
$g$ for every index. A finite signed measure $\mu$ on $\Omega\times X_{\mathcal B}^S$,
where the finite set $\Omega$ records discrete data fixed by the group (maps,
component and type indices, tags), is \emph{square-affine invariant} if it is
invariant under $(\omega,z)\mapsto(\omega,gz)$ for all $g\in\mathfrak G_{\mathcal B}$.
If $e_p=1$ for every $p\in\mathcal B$, write
$\Om[\mathcal B]=\prod_{p\in\mathcal B}\Om\subseteq X_{\mathcal B}^2$ (the pairs
$(x,y)$ with $y_p-x_p\in\Qsq$ for every $p\in\mathcal B$); the \emph{product
strict-square pair law} is the uniform probability on $\Om[\mathcal B]$, which
is $\bigotimes_{p\in\mathcal B}\Unif(\Om)$.
\end{definition}

\begin{lemma}[Marginals and pair laws]\label{lem:marginals}
Let $\mathcal B$ be as in \cref{def:map:invariance}.
\begin{enumerate}[label=\textup{(\arabic*)}]
\item\label{it:marg:closure} \emph{Invariance is preserved} by the following
 operations on square-affine invariant measures: sums and real multiples;
 normalization of a nonzero nonnegative measure; passing to the variation
 $|\mu|$; multiplication by a nonnegative function $\Phi(\omega,z)$ with
 $\Phi(\omega,gz)=\Phi(\omega,z)$ (for instance the restriction to an event
 defined by the latent record $\omega$ alone, or to injectivity at every
 prime, or, when $e_p=1$, a density that is a product of functions
 $\Ap(\pm(z_s-z_{s'})_p)$); images under $(\omega,z)\mapsto(\omega',z\circ\varsigma_\omega)$,
 where $\omega'$ is a function of $\omega$ and $\varsigma_\omega\colon S'\to S$ is
 a map of index sets depending on $\omega$ (including the
 cylinder pullback of \cref{lem:slot-array}(v), forgetting coordinates);
 products $\mu_1\otimes\mu_2$ of invariant measures over disjoint prime sets
 $\mathcal B_1,\mathcal B_2$ (as a measure over $\mathcal B_1\cup\mathcal B_2$, with
 latent record $(\omega_1,\omega_2)$); and mixtures
 $\sum_{\omega_0\in\Omega_0}\kappa(\omega_0)\,\delta_{\omega_0}\otimes\mu^{\omega_0}$ on
 $(\Omega_0\times\Omega)\times X_{\mathcal B}^S$ of invariant measures
 $\mu^{\omega_0}$ indexed by an additional latent variable $\omega_0$ with law
 $\kappa$ (for example a map shared by all primes of a block, drawn before the
 field labels, which are then independent across the primes). For a prime $p$
 with $e_p=1$ and an oriented graph $G$, the uniform law on the injective
 strict-square labellings of $G$ over $\Fp$ (the law $\lambda_p(G)$ of
 \cref{def:graphs:labeling}) and the uniform law on all its strict-square
 labellings are invariant ($\mathcal B=\{p\}$).
\item\label{it:marg:one} \emph{One-slot marginals.} If $\mu$ is an invariant
 finite signed measure on $\Omega\times X_{\mathcal B}^S$, then for every
 $s\in S$ and $\omega\in\Omega$ the measure $\mu(\{\omega\}\times\{z:z_s\in\cdot\})$
 on $X_{\mathcal B}$ is $\mu(\{\omega\}\times X_{\mathcal B}^S)\,\Unif(X_{\mathcal B})$.
 In particular the one-slot marginals of $\mu$ and of $|\mu|$ are
 $\mu(1)\Unif(X_{\mathcal B})$ and $|\mu|(1)\Unif(X_{\mathcal B})$, where
 $\mu(1)$ is the total mass; and if $\mu$ is nonnegative, the law of $z_s$
 under $\mu(\cdot\mid\omega)$ is uniform for every $\omega$ of positive mass.
 These conclusions need only invariance under the translations
 $(\omega,z)\mapsto(\omega,(z_s+c)_s)$, $c\in X_{\mathcal B}$.
\item\label{it:marg:pair} \emph{Pair laws.} Suppose $e_p=1$ for every
 $p\in\mathcal B$. If $\mu$ is a nonzero nonnegative invariant measure,
 $s\ne s'$, and $\mu$-almost surely $(z_s,z_{s'})\in\Om[\mathcal B]$, then the law
 of $(z_s,z_{s'})$ under $\mu(\cdot\mid\omega)$, and under $\mu/\mu(1)$, is the
 product strict-square pair law $\Unif(\Om[\mathcal B])$.
\end{enumerate}
\end{lemma}

\begin{proof}
\ref{it:marg:closure} Each operation commutes with the action
$(\omega,z)\mapsto(\omega,gz)$. For images, $(gz)\circ\varsigma_\omega=g(z\circ\varsigma_\omega)$
because $g$ acts slotwise; for products,
$\mathfrak G_{\mathcal B_1\cup\mathcal B_2}=\mathfrak G_{\mathcal B_1}\times\mathfrak G_{\mathcal B_2}$
acts factorwise; for the variation, the image under the bijection $g$ of the
Jordan decomposition $\mu=\mu_+-\mu_-$ is again a decomposition into mutually
singular nonnegative measures, so by uniqueness $g$ preserves $\mu_\pm$ and
$|\mu|$; the other cases are immediate, injectivity being preserved because
each $g_p$ is a bijection. The uniform laws are invariant because
$z\mapsto rz+c$ ($r\in\Qsq$) is a bijection of the respective sets of
labellings: it multiplies every difference by a nonzero square and preserves
injectivity. Functions $\Ap(\pm(z_s-z_{s'})_p)$ are invariant for the same
reason.

\ref{it:marg:one} The measure $\mu(\{\omega\}\times\{z:z_s\in\cdot\})$ is
invariant under the translations $x\mapsto x+c$, $c\in X_{\mathcal B}$, which
lie in $\mathfrak G_{\mathcal B}$ and act transitively on $X_{\mathcal B}$; so it
gives the same mass to every point. Summing over $\omega$ gives the
unconditional statement, and $|\mu|$ is invariant under the same translations
by the Jordan argument of \ref{it:marg:closure}. Only the translations were
used.

\ref{it:marg:pair} $\mathfrak G_{\mathcal B}$ acts transitively on $\Om[\mathcal B]$:
given $(x,y),(x',y')\in\Om[\mathcal B]$, put
$g(w)_p=\frac{y'_p-x'_p}{y_p-x_p}(w_p-x_p)+x'_p$; the ratio is a quotient of
nonzero squares, so $g\in\mathfrak G_{\mathcal B}$, and $g(x)=x'$, $g(y)=y'$.
Conditioning on $\omega$ preserves invariance, so the law of $(z_s,z_{s'})$
under $\mu(\cdot\mid\omega)$ is an invariant probability on the finite set
$\Om[\mathcal B]$, on which $\mathfrak G_{\mathcal B}$ acts transitively; hence it is
uniform. Mixing over $\omega$ gives the statement for $\mu/\mu(1)$. Finally
$\Om[\mathcal B]=\prod_p\Om$, so its uniform law is the product of the uniform
laws of the factors.
\end{proof}

\subsection{Independent atoms: product laws and $\alpha^{72}\le Y$}\label{subsec:pos:products}

The reference law of the functional is a product of laws of independent
atoms, finite sets of primes. The next lemma collects what the product
inherits from its factors.

An \emph{atom} is a finite set $\Blk$ of odd primes with exponents $e_p\ge1$;
$X_\Blk=\prod_{p\in\Blk}\Z/p^{e_p}\Z$ carries its uniform probability measure. Let
$V$ be a slot set with a family $\mathfrak F$ of cuts that contains a sorting
sequence, $d=|V|$, and let $\mathcal A$ be a set of pairwise disjoint atoms,
each with a measure $\mu_\Blk$ on $X_\Blk^V$. For a finite set
$\mathcal S\subseteq\mathcal A$ put $X_{\mathcal S}=\prod_{\Blk\in\mathcal S}X_\Blk$,
identify a family $(z_{v,\Blk})_{v\in V,\Blk\in\mathcal S}$ with $(z_v)_{v\in V}$,
$z_v=(z_{v,\Blk})_\Blk\in X_{\mathcal S}$, and let
\[
 \mu_{\mathcal S}=\bigotimes_{\Blk\in\mathcal S}\mu_\Blk ,\qquad
 Z_{\mathcal S}(g)=\Lambda_{\mu_{\mathcal S}}(g)\quad(g\colon X_{\mathcal S}\to\R).
\]
For $\mathcal S_1\subseteq\mathcal S$, $\E_{\mathcal S_1}g$ is the uniform average
of $g$ over the coordinates of the atoms of $\mathcal S_1$:
$\E_{\mathcal S_1}g(x)=\E_{y\in X_{\mathcal S_1}}g(x+y)$, with $y$ extended by $0$
at the other atoms.

\begin{lemma}[Signed H\"older inequality, mean and averaging]\label{lem:record:holder}
Suppose that every $\mu_\Blk$ is a probability measure on $X_\Blk^V$, reflection
positive at every cut of $\mathfrak F$, and square-affine invariant
(\cref{def:map:invariance}, for the primes of $\Blk$). Let
$\mathcal S\subseteq\mathcal A$ be finite and let $g,g_v$ ($v\in V$) be real
functions on $X_{\mathcal S}$.
\begin{enumerate}[label=(\alph*)]
\item $\mu_{\mathcal S}$ is a probability measure, RP at every cut of
 $\mathfrak F$. Hence $Z_{\mathcal S}(g)\ge0$,
 $|\Lambda_{\mu_{\mathcal S}}((g_v)_v)|\le\prod_vZ_{\mathcal S}(g_v)^{1/d}$, and
 $\|g\|_{\mathcal S}=Z_{\mathcal S}(g)^{1/d}$ is a seminorm.
\item $\mu_{\mathcal S}$ is invariant under
 $\mathfrak G_{\mathcal S}=\prod_{\Blk\in\mathcal S}\mathfrak G_\Blk$; its one-slot
 marginals are $\Unif(X_{\mathcal S})$; and for $u,w\in V$ the law of
 $(z_u,z_w)$ is the product over $\Blk\in\mathcal S$ of the laws of
 $(z_{u,\Blk},z_{w,\Blk})$ under $\mu_\Blk$. Consequently
 $|\E_xg(x)|\le\|g\|_{\mathcal S}$.
\item $\|g\circ\psi\|_{\mathcal S}=\|g\|_{\mathcal S}$ for every primewise map
 $\psi\in\mathfrak G_{\mathcal S}$, $\psi(z)_p=r_pz_p+c_p$.
\item If $\mathcal S_1\subseteq\mathcal S$, then
 $Z_{\mathcal S}(\E_{\mathcal S_1}g)=Z_{\mathcal S\setminus\mathcal S_1}(\E_{\mathcal S_1}g)
 \le Z_{\mathcal S}(g)$. If $g$ depends only on the coordinates of the atoms of
 $\mathcal S_2\subseteq\mathcal S$, then $Z_{\mathcal S}(g)=Z_{\mathcal S_2}(g)$.
\item If $\E_xg(x)=\alpha\ge0$, then $\alpha^d\le Z_{\mathcal S}(g)$. In
 particular, for the lift, the inputs and the functional of
 \cref{def:record:functional} (with $\mathcal S=\mathcal A_N$, $d=72$, and $I,D_0,s$ as there),
 \[
  \E_{z\in G_N}g^{I,D_0,s}(z)=\frac{D_0\,\#\{n\in A\cap I:\ n\equiv s\ (D_0)\}}{|I|},
 \]
 so $\E F_A=\alpha:=|A|/N$ and $\alpha^{72}\le\Yf(\Lsc,A)$.
\end{enumerate}
\end{lemma}

\begin{proof}
(a) A product of probability measures is a probability measure. At a cut
$\vartheta\in\mathfrak F$, $C_\vartheta(\mu_{\mathcal S})=\bigotimes_\Blk C_\vartheta(\mu_\Blk)$
by \eqref{eq:cc:kron}, a Kronecker product of symmetric positive semidefinite
matrices. Since $\mathfrak F$ contains a sorting sequence, \cref{thm:holder}
and \cref{cor:seminorm} apply; $d$ is even because $V$ has a cut.
(b) The invariance is the product clause of
\cref{lem:marginals}\ref{it:marg:closure}, and the one-slot marginals are
uniform by \cref{lem:marginals}\ref{it:marg:one}. Under a product measure the
coordinates of different factors are independent, and $(z_u,z_w)$ is the
family of the pairs $(z_{u,\Blk},z_{w,\Blk})$. The mean bound is the second
part of \cref{cor:seminorm}.
(c) By (b), the map $(z_v)_v\mapsto(\psi(z_v))_v$ preserves $\mu_{\mathcal S}$.
(d) $\E_{\mathcal S_1}g$ is the average of the translates $g(\cdot+y)$,
$y\in X_{\mathcal S_1}$; by (c) with $r_p=1$ they all have seminorm
$\|g\|_{\mathcal S}$, and the triangle inequality of (a) gives the inequality.
The two equalities follow from Fubini, since the factors $\mu_\Blk$ are
probability measures and the integrand does not depend on the labels of the
atoms integrated out.
(e) By (b), $\alpha\le\|g\|_{\mathcal S}$. For the lift, averaging over $z$
annihilates every nontrivial character, and the trivial character
has conductor $1$, at most the cutoff $Q$ of the lift
(\cref{sec:lifts}), so it is retained and the mean of
$\mathcal L_{I,N}f$ is
$\widehat f_I(0)=|I|^{-1}\sum_{n\in I}f(n)$. With
$f=(D_0\one_{n\equiv s\,(D_0)}\one_A)|_I$ this is the stated number; with
$I=[N]$ and $D_0=1$ it is $\alpha$. Apply the first part to $g=F_A$, for which
$Z_{\mathcal A_N}(F_A)=\Yf(\Lsc,A)$.
\end{proof}

For the slot array and its family $\mathfrak F$ the hypotheses of the lemma
are contained in \ref{it:record:R1}--\ref{it:record:R3} of
\cref{hyp:record:reference}, which Part~V verifies atom by atom at every cut of
$\mathfrak F$, not only at the sorting cuts (see the discussion after
\cref{hyp:record:reference}).

\paragraph{Consequences and what to read next.}
Only linear algebra (the Loewner order of \cref{def:prelim:operators},
Kronecker products, polar decomposition, Cauchy--Schwarz and Jensen's
inequality) and the arithmetic notation of \cref{sec:notation}: the nonzero
squares $\Qsq$, the strict-square pairs $\Om$ and the indicator $\Ap$. Later
definitions are named only as examples or applications: the laws
$\lambda_p(G)$ (\cref{def:graphs:labeling}) in \cref{lem:marginals}, the
coupled index (\cref{def:hyb:coupled}) in \cref{rem:cc:limits}, and the lift,
the functional and the reference hypotheses
(\cref{def:record:functional,hyp:record:reference}) in and after
\cref{lem:record:holder}(e). No property of square-difference-free sets is
used in this section.
